% !TeX program = xelatex
\documentclass[UTF8,a4paper,11pt,fontset=fandol]{ctexart}
\usepackage[margin=25mm]{geometry}
\usepackage{amsmath,amssymb}
\pagestyle{plain}
\setlength{\parindent}{0pt}
\setlength{\parskip}{0.55em}
\newcommand{\ind}{\mathbf{1}}

\begin{document}
\begin{center}
  {\Large\bfseries STRank+BCE 原方法：单细胞基因检出与计数排序}
\end{center}

\textbf{数据。}
同一乳腺癌标本的两张 Xenium 切片 Rep1、Rep2 轮流作为源片和测试片。
源片按空间位置划分训练集 \(T\) 和验证集 \(V\)。每个细胞提供一对观测
\[
  (x_i,y_i)\in\mathbb{R}^{1024}\times\mathbb{N}_0^G,\qquad G=313 .
\]
\(x_i\) 是冻结的 UNI 编码器提取的 H\&E 图像特征，
\(y_{ig}\) 是基因 \(g\) 在细胞 \(i\) 中的原始转录本计数。

\textbf{模型。}
用 \(T\) 的逐维均值和标准差将 \(x_i\) 标准化为 \(\tilde x_i\)。
单隐层 MLP 共享一个 256 维表示，接两个线性输出头。
\[
  h_i=\operatorname{Dropout}_{0.1}\!\bigl(\operatorname{GELU}(W\tilde x_i+b)\bigr),
  \qquad r_i=W_rh_i+b_r,\quad a_i=W_ah_i+b_a .
\]
计数头预测同一细胞内各基因的计数比例 \(p_{ig}\)，检出头估计检出概率 \(q_{ig}\)。
\[
  p_{ig}=\frac{\exp r_{ig}}{\sum_{k=1}^{G}\exp r_{ik}},\qquad
  q_{ig}=\frac{1}{1+\exp(-a_{ig})}.
\]

\textbf{监督目标。}
记 \(n_i=\sum_g y_{ig}\) 为细胞 \(i\) 的面板总计数，\(\bar n_T=|T|^{-1}\sum_{i\in T}n_i\)。
给定 \(n_i\)，计数头用多项分布描述这些计数在基因间的分配。检出头将单计数视为检出，
以 \(z_{ig}=\ind\{y_{ig}>0\}\) 为二元标签。两个头分别建模
\[
 y_i\mid x_i,n_i\sim\operatorname{Multinomial}(n_i,p_i),\qquad
 z_{ig}\mid x_i\sim\operatorname{Bernoulli}(q_{ig}).
\]
以共享层和两个输出头的权重与偏置为参数 \(\theta\)，最小化两项损失之和。
\[
\begin{aligned}
  &\min_{\theta}\quad \mathcal L_{\rm count}(\theta)+\mathcal L_{\rm det}(\theta),\\[2pt]
  \mathcal L_{\rm count}
    &= -\frac{1}{|T|\,\bar n_T}\sum_{i\in T}\sum_{g=1}^{G} y_{ig}\log p_{ig},\\
  \mathcal L_{\rm det}
    &= -\frac{1}{|T|\,G}\sum_{i\in T}\sum_{g=1}^{G}
       \bigl[z_{ig}\log q_{ig}+(1-z_{ig})\log(1-q_{ig})\bigr].
\end{aligned}
\]
第一项是按训练集平均总计数缩放的多项负对数似然（略去常数），第二项是逐基因的二元交叉熵。
两项监督均取自同一份原始计数 \(y_i\)。

\textbf{优化与预测。}
每轮打乱训练细胞，以 2048 个细胞为一批反向传播，保持 \(\bar n_T\) 固定。
用 AdamW 更新 \(\theta\)，学习率 \(10^{-3}\)，权重衰减 \(10^{-4}\)，梯度范数裁剪至 5。
最多训练 40 轮。在 \(V\) 上逐细胞计算平均精确率（AP），以各细胞 AP 的均值选模，连续 8 轮无改善则早停。
记所选模型参数为 \(\hat\theta\)。在 \(V\) 上最大化宏平均 \(F_{1.2}\) 得到原阈值 \(\tau_0\)。
测试时选出 \(\{g:q_{ig}\ge\tau_0\}\)，再按 \(r_{ig}\) 降序排列。

\newpage
\begin{center}
  {\Large\bfseries 单计数中性复评：冻结模型，重新选择阈值}
\end{center}

\textbf{数据与决策变量。}
从源片空间验证区取 \(1{,}024\) 个细胞组成校准集 \(C\)，在另一张切片上取固定的
\(4{,}096\) 个细胞组成测试集 \(E\)。两组均有相同的 313 个基因。
冻结第一页的模型参数 \(\hat\theta\)，由图像特征算出检出分数 \(q_{ig}\)。
该分数仍是模型对 \(\Pr(y_{ig}>0\mid x_i)\) 的估计。校准数据和入选规则为
\[
 \mathcal D_C=\{(q_{ig},y_{ig}):i\in C,\;1\le g\le G\},\qquad
 s_{ig}(\tau)=\ind\{q_{ig}\ge\tau\}.
\]
同一模型的全部细胞和基因共用一个阈值 \(\tau\)。这是复评中唯一待优化的量。

\textbf{单个决策的损失。}
将原始计数 \(y=0\) 定为阴性，\(y\ge2\) 定为阳性，\(y=1\) 定为中性。
令 \(s=1\) 表示选入、\(s=0\) 表示未选，定义
\[
 \ell(y,s)=\underbrace{\ind\{y=0\}\,s}_{\text{误选代价 }1}
          +\underbrace{2\,\ind\{y\ge2\}\,(1-s)}_{\text{漏选代价 }2}.
\]
因此 \(\ell(1,0)=\ell(1,1)=0\)。选中单计数基因仍计入集合大小，对误选和漏选的贡献均为零。
系数 2 表示漏掉一个阳性基因的代价相当于误选两个阴性基因。

\textbf{校准目标。}
对任一细胞集合 \(A\)，记 \(\mathrm{FP}_A(\tau)\)、\(\mathrm{FN}_A(\tau)\)
为全部细胞与基因上的误选、漏选总数。先对基因求和，再对细胞取平均，得到经验决策风险
\[
 R_A(\tau)=\frac{1}{|A|}\sum_{i\in A}\sum_{g=1}^{G}\ell\!\left(y_{ig},s_{ig}(\tau)\right)
          =\frac{\mathrm{FP}_A(\tau)+2\,\mathrm{FN}_A(\tau)}{|A|}.
\]
该目标衡量每个细胞的平均选集代价。固定模型后，各基因的预测概率和观测数据的似然都随之确定。
调整 \(\tau\) 只改变入选集合及其代价。

\textbf{求解与测试。}
取所有不同的校准分数作为候选阈值，另加一个高于最高分的阈值 \(t_{\varnothing}\) 表示全不选。
\[
 \mathcal T_C=\{q_{ig}:i\in C,\;1\le g\le G\}\cup\{t_{\varnothing}\},\qquad
 \hat\tau\in\operatorname*{arg\,min}_{\tau\in\mathcal T_C} R_C(\tau).
\]
将分数按降序扫描，同分的细胞与基因配对同时入选。代价并列时先选误选数较少者，再选总入选数较少者。
固定 \(\hat\tau\) 后，在目标片报告 \(R_E(\hat\tau)\)。例如，以 Rep1 为源片、随机种子为 42 时，
同一模型的原阈值为 \(0.22\)，中性校准后约为 \(0.476265\)，在 Rep2 测试集上每细胞平均代价为 \(42.5918\)。
\end{document}
